% Part: proof-theory
% Chapter: normalization
% Section: translations

\documentclass[../../../include/open-logic-section]{subfiles}

\begin{document}

\olfileid{pt}{nor}{tr}

\olsection{স্বাভাবিক \Log{N2i}-\usetoken{P}{proof} ও \Log{G2i}-এর মধ্যে অনুবাদ}

\Log{N2i}-র !!^{proof}s-কে \Log{G2i}-র
!!{proof}s-এ অনুবাদ করা যায়, এবং উল্টোটিও যায়।
প্রথমে লক্ষ করি, কর্তনহীন \Log{G2i}
!!{proof}s-এর অনুবাদে স্বাভাবিক
\Log{N2i}-!!{proof}s পাওয়া যায়।
% BN-SRC-584 TARGET-CORRECTION-BEGIN
\par\smallskip\noindent\textnormal{\small\textbf{উৎস-সংশোধন (BN-SRC-584):} প্রথম অনুচ্ছেদে উৎসের \Log{2i} কোনো নির্দিষ্ট প্রমাণপদ্ধতি নয়; শিরোনাম ও পরের প্রস্তাবনা অনুযায়ী কর্তনহীন \Log{G2i} বোঝানো হয়েছে।}
% BN-SRC-584 TARGET-CORRECTION-END

ফলটি বলতে \Log{N2i} ও \Log{G2i}-এর
সিকোয়েন্টের বাঁ অংশের সম্পর্ক বোঝাতে
একটু পরিভাষা দরকার। চিহ্নযুক্ত
!!{formula}s-এর সেট~$\Gamma'$ অচিহ্নিত
!!{formula}s-এর বহুসেট~$\Gamma$-র
\emph{সঙ্গত} হবে যদি প্রতিটি
!!{formula}~$!A$-এর জন্য
$\Gamma'$-তে $x : !A$ রূপের
!!{element}s-এর সংখ্যা
$\Gamma$-তে~$!A$-এর পুনরাবৃত্তি-সংখ্যার
বেশি না হয়।

\begin{cor}
যদি $\Log{G2i} \Proves \Gamma \Sequent \Delta$,
তাহলে $\Gamma' \Sequent !C$-এর
একটি স্বাভাবিক \Log{N2i}-!!{proof}~$\delta$
আছে। এখানে $\Delta$ ফাঁকা হলে
$!C \ident \lfalse$, অন্যথায়
$\Delta=\{!C\}$; আর চিহ্নযুক্ত
!!{formula}s-এর সেট~$\Gamma'$
বহুসেট~$\Gamma$-র সঙ্গত।
অর্থাৎ প্রত্যেক !!{formula}~$!A$-এর
জন্য $\Gamma'$-তে $x : !A$ রূপের
!!{element}s-এর সংখ্যা
$\Gamma$-তে~$!A$-এর পুনরাবৃত্তি-সংখ্যার
(অর্থাৎ~$!A$-এর অনুলিপির সংখ্যার)
বেশি নয়।
% BN-SRC-585 TARGET-CORRECTION-BEGIN
\par\smallskip\noindent\textnormal{\small\textbf{উৎস-সংশোধন (BN-SRC-585):} করোলারিতে উৎসের N2c নামটি N2i করা হয়েছে; লক্ষ্যসূত্র !C-কে ভিত্তি-প্রস্তাবনা অনুযায়ী ফাঁকা বা এক-সূত্রের $\Delta$-র সঙ্গে যুক্ত করা হয়েছে।}
% BN-SRC-585 TARGET-CORRECTION-END
\end{cor}

\begin{proof}
\olref[nat][tgi]{prop:G2i-to-N2i}-এর
প্রমাণ দেখলেই ফলটি মেলে:
!!{introduction} বিধি !!{proof}s-এর শেষে
এবং !!{elimination} বিধি কেবল
!!{proof}s-এর শুরুতে যোগ করি; তাই কোনো
!!a{elimination} বিধির উপরে
!!{introduction} বিধি আসে না।
\end{proof}

তবে \olref[nat][tgi]{prop:G2i-to-N2i}-এ
\Cut{} বিধির অনুবাদে \Log{N2i}-!!{proof}~$\delta$-তে
কর্তন তৈরি হতে পারে। কোনো \Log{G2i}
!!{proof} যদি~\Cut{} দিয়ে শেষ হয়,
আর তার অব্যবহিত উপ-!!{proof}s
$\pi_1$ ও $\pi_2$ যথাক্রমে
$\Gamma_1 \Sequent !A$ এবং
$!A, \Gamma_2 \Sequent \Delta_2$-র হয়,
তাহলে~$\pi_2$-র অনুবাদ~$\delta_2$-তে
যত $x:!A \Sequent !A$ স্বতঃসিদ্ধ
আছে, সেগুলিতে~$\pi_1$-এর অনুবাদ
$\delta_1$ কলম করে বসাই।
যদি~$\delta_1$-র শেষ বিধি
প্রধান !!{formula}~$!A$-সহ
!!a{introduction} হয় এবং~$\delta_2$-তে
$x:!A \Sequent !A$ স্বতঃসিদ্ধটি
!!a{elimination} বিধির প্রধান পূর্বধারণা
হয়, তবে~$\delta$-তে কর্তন তৈরি হবে।

\Log{N2i}-র !!{proof}s-কে~\Log{G2i}-র
!!{proof}s-এও অনুবাদ করা যায়।
\olref[nat][tni]{prop:N2-to-G2}-এর
প্রমাণে \Elim{\land}, \Elim{\lif},
\Elim{\lnot}, \Elim{\lforall}-এর
অনুবাদে ফলের \Log{G2i} !!{proof}-এ
\Cut{} প্রয়োগ আসে। তবে শুরুতে
\Log{N2i}-!!{proof}-টি স্বাভাবিক হলে
এটি এড়ানো যায়।
% BN-SRC-586 TARGET-CORRECTION-BEGIN
\par\smallskip\noindent\textnormal{\small\textbf{উৎস-সংশোধন (BN-SRC-586):} উৎসের অসম্পূর্ণ শর্তবাক্যে স্বাভাবিক হওয়ার পূর্বশর্ত নেই; পরের প্রস্তাবনার শর্ত অনুযায়ী সেটি স্পষ্ট করা হয়েছে।}
% BN-SRC-586 TARGET-CORRECTION-END

\Log{N2i}-!!{proof}~$\delta$-তে
সিকোয়েন্ট-সংঘটনের তালিকা $S_1$,
\dots,~$S_n$-কে \emph{শাখা} বলব,
যদি $S_1$ স্বতঃসিদ্ধ হয় এবং
যখনই $S_i$ কোনো প্রয়োগের
প্রধান পূর্বধারণা হয়,
$S_{i+1}$ হয় তার সিদ্ধান্ত।
অর্থাৎ, শাখা~$\delta$-তে স্বতঃসিদ্ধ থেকে
নিচে সিকোয়েন্ট অনুসরণ করে,
কিন্তু !!{elimination} বিধির
গৌণ পূর্বধারণায় থেমে যায়।
$S_n$ যদি শেষ-সিকোয়েন্ট হয়,
তবে শাখাটি~$\delta$-র
\emph{প্রধান শাখা}। প্রতিটি
!!{proof}~$\delta$-তে অন্তত একটি
প্রধান শাখা থাকে, কারণ প্রতিটি
বিধির অন্তত একটি প্রধান পূর্বধারণা
আছে (শুধু \Intro{\land}-এর দুটি)।


\begin{prop}\ollabel{prop:normal-N2i-to-G2i}
$\delta$ যদি~$\Gamma \Sequent !A$-এর
স্বাভাবিক $\Log{N2c}$ অথবা
$\Log{N2i}$-!!{proof} হয়,
তাহলে $\Gamma' \Sequent \Delta$-র
একটি কর্তনহীন $\Log{G2c}$-!!{proof}
(যথাক্রমে $\Log{G2i}$-!!{proof})~$\pi$
আছে। এখানে $\Gamma$ থেকে চিহ্নগুলি
সরিয়ে পাওয়া !!{formula}s-এর বহুসেট
হলো~$\Gamma'$; $!A$ যদি~$\lfalse$ না হয়
তবে $\Delta = \{!A\}$, আর তা হলে
$\Delta = \emptyset$।
\end{prop}

\begin{proof}
এবার~$\Gamma \Sequent !A$-এর
!!{proof}~$\delta$-তে বিধিপ্রয়োগের
সংখ্যার উপর আরোহ করি।
$\delta$-তে কোনো প্রয়োগ না থাকলে
সেটি $x:!A \Sequent !A$ স্বতঃসিদ্ধ;
তখন~$\pi$ সংশ্লিষ্ট
$!A \Sequent !A$ স্বতঃসিদ্ধ,
অথবা $!A \ident \lfalse$ হলে
$\lfalse \Sequent \ $ স্বতঃসিদ্ধ।

$\delta$ কোনো বিধিপ্রয়োগে শেষ হলে
ক্ষেত্রভেদ করি:
\begin{enumerate}
  \item $R$ যদি !!a{introduction}
  বিধি বা~\FalseInt হয়,
  \olref[nat][tni]{prop:N2-to-G2}-এর
  প্রমাণের মতোই অনুবাদ হয়;
  \Cut{} লাগে না।

  $\delta$-র শেষ বিধি~$R$
  যদি !!a{elimination} হয়,
  নিচের বিষয়গুলি লক্ষ করি:
  \begin{enumerate}
    \item $\delta$-র প্রধান শাখায়
    কোনো !!{introduction} বিধি নেই।
    থাকলে সেই শাখায়
    !!a{introduction} বিধি বা~\FalseInt{}-এর
    পরে !!a{elimination} বিধি থাকত,
    ফলে $\delta$ স্বাভাবিক হতো না।
    \item $\delta$-র প্রধান শাখায়
    !!{introduction} বিধি নেই বলে
    প্রধান শাখা একটিই।
    (শুধু \Intro{\land}-এর দুটি
    প্রধান পূর্বধারণা; একাধিক প্রধান
    শাখা হলে ওই বিধির পূর্বধারণা
    দিয়ে যেতে হতো।)
    \item প্রধান শাখায়
    \Elim{\lor} বা \Elim{\lexists}-এর
    প্রধান পূর্বধারণা থাকলে
    এমন বিধি মাত্র একটি এবং
    সেটিই শেষ বিধি~$R$।
    (অন্যথায় প্রধান শাখায়
    ওই বিধির সিদ্ধান্ত পরের
    !!a{elimination} বিধির
    প্রধান পূর্বধারণা হতো;
    বিন্যাস-রূপান্তর সম্ভব বলে
    প্রমাণটি স্বাভাবিক থাকত না।)
  \item \Elim{\lor} ও
  \Elim{\lexists} ছাড়া অন্য
  !!{elimination} বিধি অনুমান
  অবমুক্ত করে না, অর্থাৎ
  সিকোয়েন্টের বাঁ প্রসঙ্গ বদলায় না।
  ওই দুই বিধিও কেবল গৌণ
  পূর্বধারণার উপরের অনুমান
  অবমুক্ত করে; তাই শুধু ওই
  পূর্বধারণার বাঁ প্রসঙ্গ বদলায়।
  অর্থাৎ, $\delta$-র প্রধান শাখায়
  $S_i$-র প্রসঙ্গ~$\Gamma_1$ হলে,
  যে প্রয়োগে $S_i$ প্রধান পূর্বধারণা
  তার সিদ্ধান্ত~$S_{i+1}$-এর প্রসঙ্গ
  $\Gamma_1' \supseteq \Gamma_1$।
  \item সুতরাং প্রধান শাখার
  সর্বোচ্চ সিকোয়েন্ট~$S_1$
  যদি $x: !C \Sequent !C$ হয়,
  তবে $x:!C \in \Gamma$।
  \end{enumerate}
  এবার~$!C$-এর গঠন অনুযায়ী
  ক্ষেত্রভেদ করি। প্রধান শাখার
  পরবর্তী প্রতিটি ধাপে $S_i$
  !!a{elimination} বিধির প্রধান
  পূর্বধারণা; তাই $x:!C \Sequent !C$-এর
  ক্ষেত্রেও তা সত্য।

  \item $!C \ident !D \land !E$।
  তখন $\delta$-র রূপ নিচেরটি:
  \[
  \Axiom$x:!D \land !E \fCenter !D \land !E$
  \RightLabel{\Elim{\land}[1]}
  \UnaryInf$x:!D \land !E \fCenter !D$
  \Deduce$x:!D \land !E, \Gamma_1 \fCenter !A$
  \DisplayProof
  \]
  প্রধান শাখায়
  $x:!D \land !E \Sequent !D$ আছে;
  সেখানে \Intro{\lif} বা
  \Intro{\lnot}-এর প্রধান পূর্বধারণা,
  কিংবা \Elim{\lor} বা
  \Elim{\lexists}-এর গৌণ পূর্বধারণা
  নেই। তাই $x:!D \land !E$-এর
  প্রসঙ্গ-!!{formula}s প্রধান শাখা
  ধরে অপরিবর্তিত থাকে। এই জন্যই
  শেষ-সিকোয়েন্টের প্রসঙ্গে
  $x:!D \land !E$ থাকতে হবে।
  আরও, \Elim{\land}-এ শেষ হওয়া
  উপ-!!{proof}-এর বদলে
  $x:!D \Sequent !D$ বসিয়ে
  এবং প্রধান শাখার প্রতিটি
  সিকোয়েন্টের প্রসঙ্গে
  $x:!D \land !E$-র জায়গায়
  $x:!D$ বসিয়ে !!{proof}~$\delta$-কে
  !!{proof}~$\delta_1$-এ
  রূপ দিতে পারি; অর্থাৎ
  \[
  \bottomAlignProof
  \Axiom$x:!D \land !E \fCenter !D \land !E$
  \RightLabel{\Elim{\land}[1]}
  \UnaryInf$x:!D \land !E \fCenter !D$
  \Deduce$x:!D \land !E, \Gamma_1 \fCenter !A$
  \DisplayProof
  \quad\text{এর বদলে}\quad
  \bottomAlignProof
  \Axiom$x:!D \fCenter !D$
  \Deduce$x:!D, \Gamma_1 \fCenter !A$
  \DisplayProof
  \]
  $\delta_1$-এ আরোহ-অনুমান
  প্রযোজ্য, কারণ $\delta$-তে
  বিধিপ্রয়োগের সংখ্যা
  $\delta_1$-এর চেয়ে~$1$ বেশি।
  % BN-SRC-587 TARGET-CORRECTION-BEGIN
  \par\smallskip\noindent\textnormal{\small\textbf{উৎস-সংশোধন (BN-SRC-587):} আরোহে নির্মিত উপপ্রমাণটি $\delta_1$; উৎসের পরের গণনাবাক্যে অনির্ধারিত $\delta'$ ছিল।}
  % BN-SRC-587 TARGET-CORRECTION-END

  আরোহ-অনুমান থেকে
  $!D, \Gamma_1' \Sequent !A$-র
  একটি \Log{G2i}-!!{proof}~$\pi_1$
  পাই। \LeftR{\land} প্রয়োগ করে
  !!{proof}~$\pi$ পাই:
  % BN-SRC-588 TARGET-CORRECTION-BEGIN
  \par\smallskip\noindent\textnormal{\small\textbf{উৎস-সংশোধন (BN-SRC-588):} উৎসগদ্যে $\pi_1$-এর G2i পূর্বধারণায় x-চিহ্ন ও অবদলিত $\Gamma_1$ আছে; প্রদর্শিত G2i বৃক্ষ অনুযায়ী চিহ্নহীন $\Gamma_1'$ নেওয়া হয়েছে।}
  % BN-SRC-588 TARGET-CORRECTION-END
  \[
  \AxiomC{}
  \RightLabel{$\pi_1$}
  \Deduce$!D, \Gamma_1' \fCenter !A$
  \RightLabel{\LeftR{\land}}
  \UnaryInf$!D \land !E, \Gamma_1' \fCenter !A$
  \DisplayProof
  \]
  $!A \ident \lfalse$ হলে
  একটি \RightR{\Weakening}
  বসাই, যেমন:
  \[
  \AxiomC{}
  \RightLabel{$\pi_1$}
  \Deduce$!D, \Gamma_1' \fCenter $
  \RightLabel{\LeftR{\land}}
  \UnaryInf$!D \land !E, \Gamma_1' \fCenter $
  \RightLabel{\RightR{\Weakening}}
  \UnaryInf$!D \land !E, \Gamma_1' \fCenter !A$
  \DisplayProof
  \]
  \item $!C \ident !D \lor !E$ হলে
  সেটি \Elim{\lor}-এর প্রধান
  পূর্বধারণা হতে হবে, এবং এই
  প্রয়োগটিই প্রধান শাখার
  শেষ প্রয়োগ~$R$। অর্থাৎ,
  $\delta$-র রূপ নিচেরটি:
  \[
  \Axiom$x:!D \lor !E \fCenter !D \lor !E$
  \AxiomC{}
  \RightLabel{$\delta_1$}
  \Deduce$y:!D, \Gamma_1 \fCenter!A$
  \AxiomC{}
  \RightLabel{$\delta_2$}
  \Deduce$z:!E, \Gamma_2 \fCenter !A$
  \RightLabel{\Elim{\lor}}
  \TrinaryInf$x:!D \lor !E, \Gamma_1, \Gamma_2 \fCenter !A$
  \DisplayProof
  \]
  !!{proof}s $\delta_1$ ও
  $\delta_2$-তে আরোহ-অনুমান
  প্রয়োগ করে যথাক্রমে
  $!D, \Gamma_1' \Sequent !A$
  ও $!E, \Gamma_2' \Sequent !A$-র
  \Log{G2i}-!!{proof}s
  $\pi_1$ ও $\pi_2$ পাই।
  $\LeftR{\lor}$ প্রয়োগে~$\pi$ পাই:
  \[
  \AxiomC{}
  \RightLabel{$\pi_1$}
  \Deduce$!D, \Gamma_1' \fCenter !A$
  \AxiomC{}
  \RightLabel{$\pi_2$}
  \Deduce$!E, \Gamma_2' \fCenter !A$
  \RightLabel{\LeftR{\lor}}
  \BinaryInf$!D \lor !E, \Gamma_1', \Gamma_2' \fCenter !A$
  \DisplayProof
  \]
  $\Elim{\lor}$ যদি
  $y: !D$ বা $z: !E$ অবমুক্ত না করে
  (অর্থাৎ সেগুলি গৌণ
  পূর্বধারণাগুলির প্রসঙ্গে
  না থাকে), অথবা~$!A$
  যদি $\lfalse$ হয়, তবে
  কয়েকটি \LeftR{\Weakening}
  ও \RightR{\Weakening}
  যোগ করতে হয়। যেমন:
  \[
  \AxiomC{}
  \RightLabel{$\pi_1$}
  \Deduce$\Gamma_1' \fCenter $
  \RightLabel{\LeftR{\Weakening}}
  \UnaryInf$!D, \Gamma_1' \fCenter $
  \AxiomC{}
  \RightLabel{$\pi_2$}
  \Deduce$\Gamma_2' \fCenter $
  \RightLabel{\LeftR{\Weakening}}
  \UnaryInf$!E, \Gamma_2' \fCenter $
  \RightLabel{\LeftR{\lor}}
  \BinaryInf$!D \lor !E, \Gamma_1', \Gamma_2' \fCenter $
  \RightLabel{\RightR{\Weakening}}
  \UnaryInf$!D \lor !E, \Gamma_1', \Gamma_2' \fCenter !A$
  \DisplayProof
  \]

  \item $!C \ident !D \lif !E$।
  তখন $\delta$-র রূপ নিচেরটি:
  \[
  \Axiom$x:!D \lif !E \fCenter !D \lif !E$
  \AxiomC{}
  \RightLabel{$\delta_1$}
  \Deduce$\Gamma_1 \fCenter!D$
  \RightLabel{\Elim{\lif}}
  \BinaryInf$x:!D \lif !E, \Gamma_1 \fCenter !E$
  \RightLabel{$\delta_2$}
  \Deduce$x:!D \lif !E, \Gamma_1, \Gamma_2 \fCenter !A$
  \DisplayProof
  \]
  লক্ষ করুন, প্রধান শাখায়
  $x:!D \lif !E, \Gamma_1 \Sequent !E$
  থাকলেও তাতে \Intro{\lif}
  বা \Intro{\lnot}-এর
  প্রধান পূর্বধারণা কিংবা
  \Elim{\lor} বা \Elim{\lexists}-এর
  গৌণ পূর্বধারণা নেই।
  তাই $x:!D \lif !E, \Gamma_1$-এর
  প্রসঙ্গ-!!{formula}s প্রধান শাখায়
  অপরিবর্তিত থাকে। এই জন্য
  শেষ-সিকোয়েন্টের প্রসঙ্গে
  $x:!D \lif !E, \Gamma_1$
  থাকতে হবে। আরও,
  \Elim{\lif}-এ শেষ হওয়া
  উপ-!!{proof}-এর বদলে
  $x:!E \Sequent !E$ বসিয়ে,
  এবং প্রধান শাখার প্রতিটি
  সিকোয়েন্টের প্রসঙ্গে
  $x:!D \lif !E, \Gamma_1$-র
  জায়গায় $x:!E$ বসিয়ে
  !!{proof}-খণ্ড~$\delta_2$-কে
  সঠিক !!{proof}~$\delta_2'$-তে
  রূপ দিতে পারি; অর্থাৎ
  \[
  \bottomAlignProof
  \Axiom$x:!D \lif !E, \Gamma_1 \fCenter !E$
  \RightLabel{$\delta_2$}
  \Deduce$x:!D \lif !E, \Gamma_1, \Gamma_2 \fCenter !A$
  \DisplayProof
  \quad\text{এর বদলে}\quad
  \bottomAlignProof
  \Axiom$x:!E \fCenter !E$
  \RightLabel{$\delta_2'$}
  \Deduce$x:!E, \Gamma_2 \fCenter !A$
  \DisplayProof
  \]
  $\delta_1$ ও $\delta_2'$-তে
  আরোহ-অনুমান প্রযোজ্য:
  $\delta$-তে বিধিপ্রয়োগের সংখ্যা
  $\delta_1$ ও~$\delta_2$-তে
  প্রয়োগের মোট সংখ্যার
  চেয়ে~$1$ বেশি; এই অতিরিক্তটি
  প্রধান শাখার শুরুতে থাকা
  \Elim{\lif} প্রয়োগ।
  (ওই প্রয়োগই $\delta$-র
  শেষ প্রয়োগ~$R$ হলে
  $\delta_1$-তে $0$টি
  প্রয়োগ থাকে।)

  আরোহ-অনুমান থেকে
  $\Gamma_1' \Sequent !D$
  ও $!E, \Gamma_2' \Sequent !A$-র
  \Log{G2i}-!!{proof}s
  যথাক্রমে $\pi_1$ ও $\pi_2$ পাই।
  \LeftR{\lif} প্রয়োগে
  !!{proof}~$\pi$ পাই:
  % BN-SRC-589 TARGET-CORRECTION-BEGIN
  \par\smallskip\noindent\textnormal{\small\textbf{উৎস-সংশোধন (BN-SRC-589):} উৎসগদ্যে G2i পূর্বধারণাদুটির প্রসঙ্গ যথাক্রমে $\Gamma_1$ ও $\Gamma_1'$; প্রদর্শিত বিধিবৃক্ষ এবং চিহ্ন-অপসারণ অনুযায়ী সেগুলি $\Gamma_1'$ ও $\Gamma_2'$।}
  % BN-SRC-589 TARGET-CORRECTION-END
  \[
  \AxiomC{}
  \RightLabel{$\pi_1$}
  \Deduce$\Gamma_1' \fCenter !D$
  \AxiomC{}
  \RightLabel{$\pi_2$}
  \Deduce$!E, \Gamma_2' \fCenter !A$
  \RightLabel{\LeftR{\lif}}
  \BinaryInf$!D \lif !E, \Gamma_1', \Gamma_2' \fCenter !A$
  \DisplayProof
  \]
  $!D$ এবং/অথবা
  $!A \ident \lfalse$ হলে
  একটি বা দুটি
  \RightR{\Weakening}
  বসাই, যেমন:
  \[
  \AxiomC{}
  \RightLabel{$\pi_1$}
  \Deduce$\Gamma_1' \fCenter $
  \RightLabel{\RightR{\Weakening}}
  \UnaryInf$\Gamma_1' \fCenter !D$
  \AxiomC{}
  \RightLabel{$\pi_2$}
  \Deduce$!E, \Gamma_2' \fCenter $
  \RightLabel{\RightR{\Weakening}}
  \UnaryInf$!E, \Gamma_2' \fCenter !A$
  \RightLabel{\LeftR{\lif}}
  \BinaryInf$!D \lif !E, \Gamma_1', \Gamma_2' \fCenter !A$
  \DisplayProof
  \]

\end{enumerate}
অবশিষ্ট ক্ষেত্রগুলিও একইভাবে
বিচার করতে হয়; সেগুলি
অনুশীলন হিসেবে রাখা হলো।
\end{proof}

\begin{cor}
স্বাভাবিক \Log{N1i} অথবা
\Log{N2i}-!!{proof}-এ থাকা
প্রতিটি !!{formula} হলো
শেষ-!!{formula} বা
শেষ-সিকোয়েন্টে থাকা কোনো
সূত্রের উপ-!!{formula};
\Log{N1i}-র ক্ষেত্রে
অবমুক্ত না হওয়া অনুমানের
সূত্রগুলির উপসূত্রও হতে পারে।
% BN-SRC-590 TARGET-CORRECTION-BEGIN
\par\smallskip\noindent\textnormal{\small\textbf{উৎস-সংশোধন (BN-SRC-590):} N1i-র খোলা অনুমানের সূত্রও স্বাভাবিক প্রমাণে থাকতে পারে; শুধু শেষসূত্রের উপসূত্র বললে সেই ক্ষেত্র বাদ যায়। আগের উৎস-সংশোধনের একই পূর্বধারণা-সহ উপসূত্র-ধর্ম এখানে স্পষ্ট হয়েছে।}
% BN-SRC-590 TARGET-CORRECTION-END
\end{cor}

\begin{proof}
  \olref{prop:normal-N2i-to-G2i}
  অনুযায়ী প্রতিটি স্বাভাবিক
  \Log{N2i}-!!{proof}-কে
  \Log{G2i}-!!{proof}-এ
  অনুবাদ করা যায়। প্রমাণটি
  দেখলে বোঝা যায়,
  \Log{N2i}-!!{proof}-এ থাকা
  প্রতিটি !!{formula}
  \Log{G2i}-!!{proof}-এও
  থাকে। \Log{G2i}-তে
  \Cut{} বিধি নেই; তাই
  তার উপ-!!{formula}-ধর্ম আছে।
\end{proof}
% BN-NORM-174 TARGET-NORMALIZATION-BEGIN
\par\smallskip\noindent\textnormal{\small\textbf{ভাষা-স্বাভাবিকীকরণ (BN-NORM-174):} দুটি প্রদর্শিত রূপান্তরের মাঝের ইংরেজি into সংযোগপদ বাংলা এর বদলে হয়েছে; সূত্র, বিধি ও শাখা অপরিবর্তিত।}
% BN-NORM-174 TARGET-NORMALIZATION-END

\end{document}
